% ============================================================
%  Função Modular — Apostila Premium | PratiqueJá
%  Engine: XeLaTeX
% ============================================================
\documentclass[12pt]{article}
\usepackage{../../pratiqueja}
\usepackage{tikz}

% \hlm — destaque de resultado em modo-math (cor sem bold)
\newcommand{\hlm}[1]{{\color{darkblue}#1}}

\setsubject{Função Modular}

\begin{document}
\pagenumbering{arabic}
\setstretch{1.2}
\color{bodytext}

\heroheader{Função Modular}%
           {Módulo, Gráfico em V, Equações e Inequações}%
           {Matemática · Avançado}

\setlength{\columnsep}{18pt}
\setlength{\columnseprule}{0.5pt}
\def\columnseprulecolor{\color{exborder}}

\begin{multicols}{2}

%─────────────────────────────────────────────────────────────
\TW{Def}{badgepurple}{Função Modular}{$f(x)=|x|$ — definida por casos}

\regra{O \textbf{módulo} (valor absoluto) de um real é sua \textbf{distância
à origem}: $\;|5|=5,\ \ |-3|=3,\ \ |0|=0$.}
\formula{$|x| = \begin{cases} x & \text{se } x \geq 0\\ -x & \text{se } x<0 \end{cases}$}

\obs{A \textbf{função modular} $f(x)=|g(x)|$ tem imagem $[0,+\infty)$:
resultado sempre \textbf{não negativo}.}

%─────────────────────────────────────────────────────────────
\TW{Graf}{darkblue}{Gráfico — Forma em V}{Vértice e deslocamentos}

\regra{$f(x)=|x|$ tem forma de \textbf{V} com vértice na origem.
$f(x)=|x-h|+k$ desloca o vértice para $(h,k)$.}

\begin{center}
\begin{tikzpicture}[scale=0.40,font=\scriptsize,>=stealth]
\draw[->,line width=0.7pt] (-3.6,0)--(4.4,0) node[right]{$x$};
\draw[->,line width=0.7pt] (0,-0.4)--(0,4.6) node[above]{$y$};
\draw[darkblue,line width=1.2pt] (-3,3)--(0,0)--(3,3);
\draw[vered,line width=1.2pt,dashed] (-2,4)--(1,1)--(4,4);
\fill[darkblue] (0,0) circle (2.5pt);
\fill[vered] (1,1) circle (2.5pt);
\draw[graycolor,line width=0.4pt] (1.05,0.95)--(1.5,0.55);
\node[anchor=west,font=\scriptsize] at (1.5,0.5){$(1,1)$};
\node[darkblue,anchor=north] at (-2.0,1.0){$|x|$};
\node[vered,anchor=west] at (0.25,3.9){$|x{-}1|{+}1$};
\end{tikzpicture}
\end{center}

\obs{\textbf{Dom:} $\mathbb{R}$; \textbf{Im:} $[k,+\infty)$. Decrescente
em $(-\infty,h]$, crescente em $[h,+\infty)$.}

%─────────────────────────────────────────────────────────────
\TW{Eq}{badgegreen}{Equações Modulares}{$|f(x)|=k$ — dois casos}

\regra{$|f(x)| = k$ $(k\geq0) \Rightarrow f(x)=k$ ou $f(x)=-k$.
Se $k<0$: $S=\emptyset$.}

\exemp{%
  \textbf{$|2x-3|=5$:}\\
  $2x-3=5 \Rightarrow x=4$; \ $2x-3=-5 \Rightarrow x=-1$\\
  \hlm{$S=\{-1,\,4\}$}%
}
\regra{\textbf{Dois módulos} — $|f(x)|=|g(x)| \Rightarrow f=g$ ou $f=-g$.}
\exemp{%
  \textbf{$|x+2|=|x-1|$:}\\
  $x+2=x-1$ (impossível); \ $x+2=-(x-1) \Rightarrow x=-\tfrac12$\\
  \hlm{$S=\left\{-\tfrac12\right\}$}%
}

%─────────────────────────────────────────────────────────────
\TW{Ineq}{badgepurple}{Inequações Modulares}{Intervalo ou união}

\regra{Com $k>0$:\\[2pt]
\textbf{$|f(x)|<k$} → intervalo $-k<f(x)<k$ \ {\footnotesize(\emph{menor: and})}\\[2pt]
\textbf{$|f(x)|>k$} → união $f(x)<-k$ ou $f(x)>k$ \ {\footnotesize(\emph{maior: or})}}

\exemp{%
  \textbf{$|x+1|\leq4$:} \ $-4\leq x+1\leq4 \Rightarrow \hlm{S=[-5,3]}$\\[2pt]
  \textbf{$|3x-1|>2$:} \ $x>1$ ou $x<-\tfrac13$\\
  $\Rightarrow \hlm{S=\left(-\infty,-\tfrac13\right)\cup(1,+\infty)}$%
}

\end{multicols}

\end{document}
